\section*{Facilities, Equipment, and Other Resources}
\addcontentsline{toc}{section}{\protect\numberline{}Facilities, Equipment, and Other Resources}

The proposed project will leverage established computational and collaborative resources at Auburn University as well as shared infrastructure provided through existing interdisciplinary research centers. These resources are sufficient to support the technical and dissemination components of the proposed research. Existing collaborative mechanisms between the participating investigators, including joint research meetings and co-mentored students, will facilitate efficient coordination across sites. These resources are sufficient to support the proposed research activities without reliance on additional external facilities.

The PIs and all participating personnel will have full access to resources provided by the Department of Computer Science and Software Engineering's Assistive Intelligence Lab directed by PI Ding at Auburn University. Auburn University is designated as a National Security Agency (NSA) and Department of Homeland Security (DHS) Center of Academic Excellence in Cyber Research (CAE-R), Cyber Operations (CAE-CO), and Cyber Defense (CAE-CD). These designations reflect strengths in cybersecurity research and education and provide an established environment for interdisciplinary collaboration, graduate training, and dissemination of research outputs. The project will draw on these capabilities to support research execution and integration of results into graduate curricula and research training activities.

\noindent\textbf{Assistive Intelligence Lab:} \\
The Assistive Intelligence Lab focuses on the development of novel AI algorithms that can better understand human behaviors to facilitate naturalistic human-AI integration. The lab has approximately 8 members including PhD, MS, Undergraduate, and visitors. Auburn's CSSE department provides space and research labs for its faculty and graduate students. the Assistive Intelligence Lab occupies approximately 500 sq ft in addition to shared meeting room and facilities spaces. The CSSE department has also public access computer labs and special purpose laboratories for research. Large conference rooms are available to practice presentations and to hold meetings. Other smaller meeting rooms are also available for collaboration and discussions. High-speed wired and wireless network access is available in all areas. The CSSE department also has a dedicated server room for use by its faculty.

\noindent\textbf{Equipment:} The PI's lab has four dedicated servers are on-site which has 8xNvidia RTX 4090, 4x RTX 5090, 4x RTX PRO 6000. One additional server with 6x RTX PRO 6000 cards for development and fine tuning of deep learning models is in the department's server room. These servers are connected to 28TB of network attached storage over a 10Gbps network. Existing students also have access to multiple laptops and 2 workstations. Access to equipment will be provided via VPN.

\noindent\textbf{Alabama Supercomputer Authority (ASA): } ASA provides High Performance Computing (HPC) services at no cost to faculty, staff and enrolled students of public schools, colleges, and universities in Alabama. This allows universities to create a trained workforce with the skills needed for jobs in engineering, science, mathematics, aerospace and medical research fields. Dr. Ding and his students have been provided access to these resources. The HPC system contains 5,056 x86-64 Proc, Shared/Distributed Memory Architecture, InfiniBand high speed/low latency network, 512GB-4TB memory per node, 33 TB total memory, 1.4 PB shared storage via Panasas, A100 and H100 GPUs. 

\noindent\textbf{Office: } All faculty members have an office, and there is adequate space to host additional graduate research assistants on this project.

\noindent\textbf{Other Resources: } Departments provides the space and basic networking services to carry out the experiments, and secretarial and administrative support as well as general-purpose office equipment, such as fax, printers, and photocopiers.

